\documentclass[french]{book}
\usepackage{teach}
\usepackage{verbatim}
\usepackage{amsmath,amsfonts,amssymb}
\usepackage{pgfplots}
%\usepackage{geometry}
%\geometry{top=1cm, bottom=1cm, left=2cm, right=2cm}
\usepackage[utf8]{inputenc}
\usepackage[french]{babel}
\redefinecolor{thm}{title}{bgcolor}{alizarin}
\redefinecolor{prop}{title}{bgcolor}{meatbrown}
\redefinecolor{defn}{title}{bgcolor}{olivine}
\definecolor{alizarin}{rgb}{0.82, 0.1, 0.26}
\definecolor{meatbrown}{rgb}{0.9, 0.72, 0.23}
\definecolor{olivine}{rgb}{0.6, 0.73, 0.45}
\usepackage{pgf,tikz,tkz-tab}

\usepackage{mathrsfs}
\usetikzlibrary{arrows}

\begin{document}


\begin{center}
\huge{ Devoir surveillé VII}
\end{center}
%\title{Devoir surveillé V}
 
\vspace{1cm}
\pagestyle{empty}

\section*{Résolution algébrique d’équation $f(x)=k$ et inéquation $f(x)>k$ }

Résoudre algébriquement l’équation $f(x)=3$ avec : \\

\begin{enumerate}
\item $f(x)=x^2 $
\item $f(x)=\dfrac{1}{x}$
\item $f(x)=\sqrt{x}$
\item $f(x)=-x+ \dfrac{2}{3} $
\end{enumerate}

Résoudre algébriquement l’inéquation $f(x)>2$ avec : \\

\begin{enumerate}
\item $f(x)=x^2 $
\item $f(x)=\dfrac{1}{x}$
\item $f(x)=\sqrt{x}$
\item $f(x)=-x+ \dfrac{2}{3} $
\end{enumerate}

\section*{ Résolution graphique d’équation $f(x)=k$ et inéquation $f(x)>k$ }

On représente la courbe de quatre fonctions notées respectivement $C_f,C_g,C_h$ et $C_k$.

\begin{enumerate}
\item Pour chaque courbe représentative, donner le nom de la fonction que la courbe représente.
\item Dresser le tableau de variation de $f(x)$
\item Résoudre graphiquement $f(x) \leq 1,5$
\item Résoudre graphiquement $k(x) = 2$
\item Résoudre graphiquement $h(x)<g(x)$
\item Résoudre graphiquement $f(x)=h(x)$
\end{enumerate}

\end{document}